% ========================================================
% BÀI TẬP TOÁN 5 - KẾT NỐI TRI THỨC VỚI CUỘC SỐNG
% Bài 20: Phép trừ số thập phân
% Định dạng: BTTeX Part III - Trắc nghiệm điền số ngắn
% ========================================================

% Câu 1
\begin{ex}
  Tính kết quả: $50 - 24{,}5$. Lấy kết quả nhân với 2 được bao nhiêu?
  \shortans{51}
  \loigiai{
    $50 - 24{,}5 = 25{,}5$. $25{,}5 \times 2 = 51$.
  }
\end{ex}

% Câu 2
\begin{ex}
  Tính kết quả: $68{,}5 - 28{,}5$.
  \shortans{40}
  \loigiai{
    $68{,}5 - 28{,}5 = 40$.
  }
\end{ex}

% Câu 3
\begin{ex}
  Tính kết quả: $100 - 37{,}5$. Lấy kết quả cộng với 0,5 được bao nhiêu?
  \shortans{63}
  \loigiai{
    $100 - 37{,}5 = 62{,}5$. $62{,}5 + 0{,}5 = 63$.
  }
\end{ex}

% Câu 4
\begin{ex}
  Tính kết quả: $85{,}75 - 35{,}75$.
  \shortans{50}
  \loigiai{
    $85{,}75 - 35{,}75 = 50$.
  }
\end{ex}

% Câu 5
\begin{ex}
  Tính giá trị: $120 - 45{,}8 - 24{,}2$.
  \shortans{50}
  \loigiai{
    $120 - (45{,}8 + 24{,}2) = 120 - 70 = 50$.
  }
\end{ex}

% Câu 6
\begin{ex}
  Tính nhanh: $45 - 12{,}3 - 7{,}7$.
  \shortans{25}
  \loigiai{
    $45 - (12{,}3 + 7{,}7) = 45 - 20 = 25$.
  }
\end{ex}

% Câu 7
\begin{ex}
  Tìm số tự nhiên $x$ biết: $x + 25{,}5 = 60{,}5$.
  \shortans{35}
  \loigiai{
    $x = 60{,}5 - 25{,}5 = 35$.
  }
\end{ex}

% Câu 8
\begin{ex}
  Tìm $x$ biết: $100 - x = 42{,}5 + 17{,}5$.
  \shortans{40}
  \loigiai{
    $42{,}5 + 17{,}5 = 60$. $x = 100 - 60 = 40$.
  }
\end{ex}

% Câu 9
\begin{ex}
  Một cuộn dây điện dài $50\text{ m}$, thợ điện đã dùng hết $23{,}5\text{ m}$. Chiều dài đoạn dây còn lại khi nhân đôi là bao nhiêu mét?
  \shortans{53}
  \loigiai{
    Đoạn còn lại là $50 - 23{,}5 = 26{,}5\text{ m}$. Nhân đôi: $26{,}5 \times 2 = 53\text{ m}$.
  }
\end{ex}

% Câu 10
\begin{ex}
  Một bao gạo có $50\text{ kg}$, người ta đã lấy ra hai lần, lần thứ nhất $14{,}5\text{ kg}$, lần thứ hai $15{,}5\text{ kg}$. Trong bao còn lại bao nhiêu ki-lô-gam gạo?
  \shortans{20}
  \loigiai{
    Trong bao còn lại: $50 - (14{,}5 + 15{,}5) = 50 - 30 = 20\text{ kg}$.
  }
\end{ex}

% Câu 11
\begin{ex}
  Một can dầu hỏa nặng $18\text{ kg}$, sau khi rót bớt dầu ra thì can dầu còn nặng $8{,}5\text{ kg}$. Lượng dầu đã rót ra nhân đôi là bao nhiêu ki-lô-gam?
  \shortans{19}
  \loigiai{
    Đã rót ra: $18 - 8{,}5 = 9{,}5\text{ kg}$. Nhân đôi được 19.
  }
\end{ex}

% Câu 12
\begin{ex}
  Tính giá trị biểu thức: $75{,}5 - (24{,}2 + 11{,}3)$.
  \shortans{40}
  \loigiai{
    $75{,}5 - 35{,}5 = 40$.
  }
\end{ex}

% Câu 13
\begin{ex}
  Một mảnh vải dài $25\text{ m}$. Người ta cắt may áo hết $8{,}5\text{ m}$ và may quần hết $6{,}5\text{ m}$. Mảnh vải còn lại bao nhiêu mét?
  \shortans{10}
  \loigiai{
    Còn lại: $25 - (8{,}5 + 6{,}5) = 25 - 15 = 10\text{ m}$.
  }
\end{ex}

% Câu 14
\begin{ex}
  Tính: $200 - 65{,}4 - 34{,}6$.
  \shortans{100}
  \loigiai{
    $200 - (65{,}4 + 34{,}6) = 200 - 100 = 100$.
  }
\end{ex}

% Câu 15
\begin{ex}
  Hai thùng chứa tất cả $80$ lít nước mắm. Thùng thứ nhất chứa $45{,}5$ lít. Thùng thứ nhất chứa nhiều hơn thùng thứ hai bao nhiêu lít?
  \shortans{11}
  \loigiai{
    Thùng hai chứa: $80 - 45{,}5 = 34{,}5$ lít. Hiệu là: $45{,}5 - 34{,}5 = 11$ lít.
  }
\end{ex}

% Câu 16
\begin{ex}
  Một ô tô chở $3$ tấn hàng, chuyến thứ nhất dỡ xuống $1{,}2$ tấn hàng, chuyến thứ hai dỡ xuống $0{,}8$ tấn hàng. Trên xe còn lại bao nhiêu tấn hàng?
  \shortans{1}
  \loigiai{
    Còn lại: $3 - (1{,}2 + 0{,}8) = 3 - 2 = 1$ tấn.
  }
\end{ex}

% Câu 17
\begin{ex}
  Tìm $x$ biết: $75 - x = 28{,}4 + 21{,}6$.
  \shortans{25}
  \loigiai{
    $28{,}4 + 21{,}6 = 50$. $x = 75 - 50 = 25$.
  }
\end{ex}

% Câu 18
\begin{ex}
  Tính giá trị biểu thức: $84{,}6 - 15{,}3 - 9{,}3$. (Sửa: $84{,}6 - 14{,}6 = 70$)
  \shortans{60}
  \loigiai{
    $84{,}6 - (15{,}3 + 9{,}3) = 84{,}6 - 24{,}6 = 60$.
  }
\end{ex}

% Câu 19
\begin{ex}
  Đường chạy dài $100\text{ m}$. Bạn Nam đã chạy được $68{,}5\text{ m}$. Nam còn phải chạy bao nhiêu mét nữa để về đích? (Lấy kết quả nhân 2).
  \shortans{63}
  \loigiai{
    Còn lại: $100 - 68{,}5 = 31{,}5\text{ m}$. Nhân 2: $31{,}5 \times 2 = 63$.
  }
\end{ex}

% Câu 20
\begin{ex}
  Một bể chứa $1\,000$ lít nước, người ta đã dùng $450{,}5$ lít vào buổi sáng và $349{,}5$ lít vào buổi chiều. Trong bể còn lại bao nhiêu lít nước?
  \shortans{200}
  \loigiai{
    Đã dùng: $450{,}5 + 349{,}5 = 800$ lít. Còn lại: $1\,000 - 800 = 200$ lít.
  }
\end{ex}

% Câu 21
\begin{ex}
  Hiệu của hai số là $25$. Nếu số bị trừ bớt đi $4{,}5$ và số trừ tăng thêm $5{,}5$ thì hiệu mới là bao nhiêu?
  \shortans{15}
  \loigiai{
    Hiệu mới là: $25 - 4{,}5 - 5{,}5 = 25 - 10 = 15$.
  }
\end{ex}

% Câu 22
\begin{ex}
  Tính nhanh: $67{,}89 - 17{,}89 - 20$.
  \shortans{30}
  \loigiai{
    $(67{,}89 - 17{,}89) - 20 = 50 - 20 = 30$.
  }
\end{ex}

% Câu 23
\begin{ex}
  Một cuộn vải dài $100\text{ m}$. Ngày đầu bán được $32{,}4\text{ m}$, ngày thứ hai bán được $27{,}6\text{ m}$. Sau hai ngày cuộn vải còn lại bao nhiêu mét?
  \shortans{40}
  \loigiai{
    Còn lại: $100 - (32{,}4 + 27{,}6) = 100 - 60 = 40\text{ m}$.
  }
\end{ex}

% Câu 24
\begin{ex}
  Khi trừ một số tự nhiên đi một số thập phân có một chữ số ở phần thập phân là $4{,}5$, một học sinh quên dấu phẩy nên lấy số tự nhiên trừ đi $45$ và được kết quả là $55$. Số tự nhiên ban đầu là bao nhiêu?
  \shortans{100}
  \loigiai{
    Số tự nhiên là: $55 + 45 = 100$.
  }
\end{ex}

% Câu 25
\begin{ex}
  Tìm $x$ biết: $x + 12{,}5 = 37{,}5$. Giá trị của $x$ là bao nhiêu?
  \shortans{25}
  \loigiai{
    $x = 37{,}5 - 12{,}5 = 25$.
  }
\end{ex}
